\subsection{MAS-11: fixed characters, central zeros and the paired determinant}
The Green construction extends to a fixed primitive Dirichlet character, and more generally to a fixed unitary Hecke character of a number field, conditional only on the same spectral-annihilator theorem. Two distinctions matter: central zeros require a finite-dimensional removal, and a character which is not self-dual gives the normalized product of its completed $L$-function and that of its functional-equation partner.

\paragraph{Completion conventions.}
Let $\chi$ be a unitary Hecke character of $C_F$, let $\chi^\vee=\chi^{-1}$, and put $Q_\chi=|d_F|N\mathfrak f_\chi$. At a real place write $\chi_v(x)=\operatorname{sgn}(x)^{\kappa_v}|x|^{it_v}$, with $\kappa_v\in\{0,1\}$ and $t_v\in\R$. At a complex place write $\chi_v(z)=(z/|z|)^{m_v}|z|_v^{it_v}$, with $m_v\in\Z$, $t_v\in\R$, and $|z|_v=|z|^2$. Define
\[
 \Gamma_{\R}(s)=\pi^{-s/2}\Gamma(s/2),\qquad
 \Gamma_{\C}(s)=2(2\pi)^{-s}\Gamma(s),
\]
\[
 \Lambda_\chi(s)=Q_\chi^{s/2}
 \prod_{v\text{ real}}\Gamma_{\R}(s+\kappa_v+it_v)
 \prod_{v\text{ complex}}\Gamma_{\C}(s+|m_v|/2+it_v)
 L_f(s,\chi).
\]
Here $L_f$ uses the primitive finite Euler factors: $(1-\chi_v(\varpi_v)q_v^{-s})^{-1}$ when unramified and $1$ when ramified. The completed functional equation is
\[
 \Lambda_\chi(s)=\epsilon_\chi\Lambda_{\chi^\vee}(1-s),
 \qquad |\epsilon_\chi|=1.
\]
These local factors and the global functional equation are the Tate conventions in \cite[\S4.7, pp.~19--22; Theorem 5.22, pp.~38--39]{poonen786}; the conductor normalization is also recorded in \cite[\S1, pp.~281--282]{rajan1998}. The gamma factor without the $2$ at each complex place differs only by a constant, which will cancel below.

If $\chi=|\cdot|^{i\tau}$, set $p_\chi(s)=(s+i\tau)(s-1+i\tau)$; otherwise set $p_\chi(s)=1$. Then
\[
 \mathcal F_\chi(s)=p_\chi(s)\Lambda_\chi(s)
\]
is entire of order one. The pole statement follows from the same Tate theorem. The entire-order-one assertion is stated in \cite[\S1, p.~282]{rajan1998} for normalized unitary characters; an arbitrary unitary character is an imaginary norm twist of a normalized one, so its completion is a vertical translate times a nonzero constant. The asserted order is therefore unchanged. Also $p_\chi(s)=p_{\chi^\vee}(1-s)$, so $\mathcal F_\chi$ satisfies the same completed functional equation. Pole cancellation introduces no zero at either simple pole. Its zeros are the zeros of the completed Hecke $L$-function, and the Euler product and functional equation place them in $0\le\Re s\le1$.

For a primitive Dirichlet character modulo $q$, one may instead use the customary constant multiple
\[
 \Lambda_\chi(s)=(q/\pi)^{(s+\kappa)/2}
 \Gamma((s+\kappa)/2)L(s,\chi),\qquad
 \kappa=(1-\chi(-1))/2,
\]
\[
 \epsilon_\chi=\frac{\tau(\chi)}{i^\kappa\sqrt q},\qquad
 \tau(\chi)=\sum_{a\bmod q}\chi(a)e^{2\pi ia/q}.
\]
See \cite[(25.15.5)--(25.15.6)]{dlmf2515} and the derivation in \cite[\S6.1; (6.2.2)]{kedlayaFE}. For $q>1$ the primitive character is nonprincipal and $p_\chi=1$; for $q=1$ use $p_\chi=s(s-1)$. Constants such as the usual factor $1/2$ in $\xi$ cancel in every normalization below.

\paragraph{The actual quotient and the central removal.}
Fix $a=1/2$, a strip half-width $\alpha>1/2$, and a weight parameter $\nu>1$. The slight enlargement of the strip ensures that all completed zeros are interior using only the closed-strip fact just established. No boundary nonvanishing theorem is needed. Put
\[
 H_\alpha=L^2(\R,e^{2\alpha|u|}(1+u^2)^{\nu/2}\,du).
\]
Choose the unitary base point $\chi$ in its compact-character component, as allowed in \cite[Notation 10.2.2.1, p.~445]{masullo}, and normalize compact Haar measure to mass one. Label by $E_\chi$ the quotient sector whose dual has base point $\chi$. In the already centered ambient realization, the unitary radial gauge
\[
 f(g)=\overline{\chi(g)}\,h(\log|g|)
\]
identifies its ambient space with $H_\alpha$ and its closed centered theta-image with some $V_\chi\subset H_\alpha$. The centering by $a=1/2$ is already incorporated in this theta-image, as in (10.2.2.1)--(10.2.2.2) of the source. This convention labels the primal sector by its dual character, eliminating any convention about conjugate isotype notation.

The sole draft input is \cite[Theorem 10.2.2.1, (10.2.2.4)--(10.2.2.6), p.~446]{masullo}, with the Hilbert-dual density described in \cite[Remark 10.2.1.2, pp.~440--441]{masullo}: on $E_\chi=H_\alpha/V_\chi$, the independent continuous jets
\[
 \ell_{\rho,j}(h)=\int_\R h(u)e^{z_\rho u}u^j\,du,
 \quad z_\rho=\rho-\tfrac12,\quad 0\le j<N_\rho,
\]
have norm-dense span in $E_\chi'$. Throughout all sums and products, $\rho$ runs over the distinct completed zeros and $N_\rho$ denotes its multiplicity. Restriction to the fixed sector follows by bounded compact-character averaging from the source's full dual density. In these coordinates pure translations are the original radial translations multiplied by the inverse unitary scalar supplied by $\chi$ on the chosen radial section. This gauge is essential for an arbitrary imaginary norm twist.

The source's $L(s,\chi)$ is already completed: \cite[Definition 3.2.1.2 and (3.2.1.2)--(3.2.1.5), p.~123]{masullo} include the archimedean factors, while Definition 3.2.1.4 separately names the finite $L_f$. Away from the polar case, our completion differs by the nowhere-zero factor $2^{r_2}Q_\chi^{s/2}$. In the polar case, Proposition 3.2.1.3(2) gives the same two simple poles, and the ideal in \cite[Theorem 7.1.4.2(1), p.~326]{masullo} uses $(s+i\tau)(1-s-i\tau)L(s,\chi)$, equal to $-p_\chi(s)L(s,\chi)$. Hence its zero set with multiplicities is exactly that of $\mathcal F_\chi$; clearing the poles contributes no extra spectral jets. The product formula (7.1.4.9) from that draft is not used here.

Put $B_0=\alpha+\nu/4$ and choose $\mu>B_0$. The convolution kernel $g_\mu(v)=e^{-\mu|v|}/(2\mu)$ induces $K_\chi$ on $E_\chi$. The proof of \fid{MAS-10} applies with precisely the following parameter replacements:
\[
 |(\log\sqrt{w_\alpha})'|\le B_0,\qquad
 F_\mu(z)=\frac1{\mu^2-z^2},\qquad
 \lambda_\rho=\frac1{\mu^2-z_\rho^2}.
\]
Indeed the complete annihilator and the translation formula preserve $V_\chi$; the weighted weak-domain calculation gives the same numerical sector with constant $B_0/\sqrt{\mu^2-B_0^2}$. Since $|\Re z_\rho|\le1/2$,
\[
 \Re\lambda_\rho>0,\qquad
 |\lambda_\rho|\le\frac1{\mu^2-1/4+(\Im\rho)^2}.
\]
Entire order one gives $\sum_{\rho\ne1/2}N_\rho|z_\rho|^{-2}<\infty$, and hence $\sum_\rho N_\rho|\lambda_\rho|<\infty$. For example, Jensen's formula gives $n(R)=O_\eta(R^{1+\eta})$ for every $\eta>0$, and dyadic summation with $\eta<1$ proves the asserted reciprocal-square sum. Gram--Schmidt on the invariant finite jet flags therefore gives the finite diagonal sum of $\Re K_\chi^*$ used in \fid{MAS-10}. Its sector-factorization lemma proves $K_\chi\in\mathcal S_1$, before any use of eigenvalue trace formulas. The same compact triangular argument identifies its exact nonzero spectrum and algebraic multiplicities. In particular,
\[
 \|K_\chi\|_1\le\frac{\mu}{\sqrt{\mu^2-B_0^2}}
 \sum_\rho N_\rho\Re\frac1{\mu^2-z_\rho^2}.
\]
Neither GRH nor the draft's nuclearity theorem is used.

Let $m=\operatorname{ord}_{s=1/2}\mathcal F_\chi(s)$, and let $P_{\chi,0}$ be the Riesz projection of $K_\chi$ for $\mu^{-2}$, or zero if $m=0$. Its rank is exactly $m$: its transpose is identity on the full central jet block and zero on every other jet block. This also proves independence of $P_{\chi,0}$ from $\mu$, by dual density. Remove its entire generalized eigenspace and set
\[
 E_\chi^\circ=(I-P_{\chi,0})E_\chi,\qquad
 K_\chi^\circ=K_\chi|_{E_\chi^\circ},\qquad
 R_\chi=K_\chi^\circ(I-\mu^2K_\chi^\circ)^{-1}.
\]
The complement need not be orthogonal. It is closed and invariant, restriction preserves trace class, and the compact spectral description makes the indicated inverse bounded. Consequently $R_\chi\in\mathcal S_1(E_\chi^\circ)$. Its complete noncentral jet multiplier is $-1/z^2$, independent of $\mu$, and its nonzero eigenvalues are $-z_\rho^{-2}$ with their full algebraic multiplicities. Coincident values are counted additively. Removing only eigenvectors need not suffice at a multiple central zero: $F_\mu$ is even at zero and can have nontrivial higher central jets.

\paragraph{The normalized determinant and the partner involution.}
Write
\[
 C_\chi(z)=\mathcal F_\chi(\tfrac12+z),\qquad
 f_\chi(z)=z^{-m}C_\chi(z),\qquad
 a_\chi=f_\chi(0)=\mathcal F_\chi^{(m)}(1/2)/m!\ne0.
\]
The functional equation gives the same central order $m$ for $\chi^\vee$, and
\[
 \frac{f_\chi(-z)}{a_\chi}
 =\frac{f_{\chi^\vee}(z)}{a_{\chi^\vee}}.
\]
Thus
\[
 \mathcal H_\chi(z)=
 \frac{f_\chi(z)f_{\chi^\vee}(z)}{a_\chi a_{\chi^\vee}}
\]
is even, entire of order at most one, and equal to $1$ at zero. Define $G_\chi(w)=\mathcal H_\chi(i\sqrt w)$ by its even power series; this is an entire function of order at most $1/2$, with no square-root branch choice. The usual conjugation identity also gives $a_{\chi^\vee}=\overline{a_\chi}$ in the displayed completion conventions.

\begin{proposition}
Under the exact spectral-annihilator input above,
\begin{equation}\label{eq:MasCharacterFull}
 \det_{E_\chi^\circ}(I-wR_\chi)=G_\chi(w).
\end{equation}
On the labeled direct sum $E_\chi^\circ\oplus E_{\chi^\vee}^\circ$, inversion induces a unitary exchange involution. The determinant on its even summand is also $G_\chi(w)$, whereas the determinant on the whole direct sum is $G_\chi(w)^2$.

If $\chi=\chi^\vee$ as a global character, inversion acts within $E_\chi^\circ$. On its even summand,
\begin{equation}\label{eq:MasCharacterSelfdual}
 \det_{E_{\chi,+}^\circ}(I-wR_{\chi,+})
 =\frac{f_\chi(i\sqrt w)}{a_\chi},
\end{equation}
and the full single-sector determinant is the square of this function.
\end{proposition}
\begin{proof}
Hadamard factorization of the central-zero-removed entire function gives
\[
 \frac{f_\chi(z)}{a_\chi}
 =e^{b_\chi z}\prod_{\rho\ne1/2}
 \left[(1-z/z_\rho)e^{z/z_\rho}\right]^{N_\rho}.
\]
Multiply by the same identity at $-z$. Both the linear exponential and every genus-one convergence factor cancel. The absolutely locally uniformly convergent result is
\[
 \frac{f_\chi(z)f_\chi(-z)}{a_\chi^2}
 =\prod_{\rho\ne1/2}(1-z^2/z_\rho^2)^{N_\rho}.
\]
The functional equation identifies the left side with $\mathcal H_\chi(z)$. The trace-class determinant theorem, with the algebraic multiplicities already proved, now gives
\[
 \det(I-wR_\chi)
 =\prod_{\rho\ne1/2}(1+w/z_\rho^2)^{N_\rho}
 =\mathcal H_\chi(i\sqrt w).
\]
This proves \eqref{eq:MasCharacterFull}, including its normalization and the absence of a residual exponential factor.

In the radial gauges, $J h(u)=h(-u)$ is a unitary map from the $\chi$ space to the $\chi^\vee$ space. The functional equation pairs the zeros $z$ and $-z$ with their multiplicities, and
\[
 \ell_{\chi^\vee,\,1-\rho,j}(Jh)=(-1)^j\ell_{\chi,\rho,j}(h).
\]
The full annihilator descriptions imply $JV_\chi=V_{\chi^\vee}$. Evenness of $g_\mu$ gives $JK_\chi=K_{\chi^\vee}J$, hence the same intertwining for the central Riesz projections and for $R_\chi$. This proves that $J$ restricts to a unitary map of the central complements. The exchange involution is
\[
 \mathcal J(x,y)=(Jy,Jx).
\]
Its even summand is the graph $\{(x,Jx)/\sqrt2:x\in E_\chi^\circ\}$, unitarily identified with $E_\chi^\circ$. The restricted direct-sum operator is therefore unitarily equivalent to $R_\chi$. The partner determinant is identical by the functional equation, proving both direct-sum assertions. Labeled copies are retained even if the compact restrictions of $\chi$ and $\chi^\vee$ coincide; their chosen unitary base points and radial gauges can differ.

When $\chi=\chi^\vee$, the functional equation gives $\epsilon_\chi=(-1)^m$ and $f_\chi(-z)=f_\chi(z)$. Reflection exchanges the two independent noncentral jet blocks at $z$ and $-z$; its even part has dimension $N_\rho$ for that pair. Applying the bounded even projector to the dense jet span proves completeness on the even quotient. The resulting determinant has precisely one factor $1+w/z_\rho^2$ per unordered pair, with full multiplicity $N_\rho$. Its square is \eqref{eq:MasCharacterFull}, while $(f_\chi(i\sqrt w)/a_\chi)^2=G_\chi(w)$. Both candidate square roots are entire and equal to $1$ at zero, so they agree near zero and hence everywhere. This proves \eqref{eq:MasCharacterSelfdual}. The central jets, before removal, have reflection signs $(-1)^j$ for $0\le j<m$; their even and odd dimensions are $\lceil m/2\rceil$ and $\lfloor m/2\rfloor$. No assumption of central nonvanishing or simple zeros enters the argument.
\end{proof}

The half-width $\alpha=1/2$ is also admissible when the classical nonvanishing theorem on $\Re s=1$ is included, because the functional equation then puts all completed zeros in the open strip. The formulation above deliberately needs only the weaker closed-strip consequence of the Euler product and functional equation. It concerns fixed characters, with no asserted bound uniform in conductor or infinity type. For a non-self-dual character, the single-sector inverse square naturally yields the paired product; an even completed factor belonging to that character alone is not supplied by its functional equation. These determinant identities establish neither positivity nor selfadjointness in the actual quotient metric.

\paragraph{Grothendieck framework.}
The conclusion is an analytic trace-class realization of completed rank-one Hecke $L$-functions on the stated adelic Hilbert quotients. Cote~1 supplies the ordinary Fredholm trace/determinant mechanism after the sector argument has proved membership in $\mathcal S_1$; its exterior-power coefficients and trace logarithm apply to $R_\chi$ and its parity restrictions. Rank one here describes the local character representations in the Euler factors. The argument supplies no constructible $\ell$-adic complex on $\Spec\mathcal O_F$, no arithmetic cohomology with a Frobenius or correspondence whose determinant is this analytic determinant, and no identification of the quotient with a motive. The geometric trace-formula comparison therefore retains the same boundary as in the preceding Grothendieck-framework discussion: the new result closes the analytic character determinant step, conditional on the spectral-annihilator theorem, without constructing such a geometric identification.

